% 25_body.tex — 산출물 ㉕ 세그먼트·JTBD·전략 커널 1페이저 (빈 템플릿 / 완성 예시 공용 본문) — gen_product_templates.py 가 Product_specs.py 에서 생성
% 드라이버: 25_세그먼트전략커널_빈템플릿.tex (\wbblanktrue) / 25_세그먼트전략커널_완성예시.tex (\wbblankfalse — 워크북 §X.5 확정 후)
\documentclass[10pt,a4paper]{article}
\usepackage[a4paper,margin=16mm,top=14mm,bottom=20mm]{geometry}
\def\DocNum{25}\def\DocDay{5}\def\DocIdx{5}
\def\DocTitle{세그먼트·JTBD·전략 커널 1페이저}
\def\DocSub{Segment, JTBD \& Strategy Kernel}
\def\DocVariant{\B{빈 템플릿}{딥게이지 완성 예시}}
\usepackage{wbstyle}
\def\DocCirc{㉕}   % newunicodechar(㉛~㊵ 폴백) 뒤에 정의해야 활성 문자로 읽힌다
\begin{document}
\docheader
 {\Bg{[회사명 · 제품명]}{딥게이지 · GAUGE}}
 {\Bg{[v0.x / D+n]}{v1.0 / 2026-05-29 (D+90)}}
 {\Bg{[작성자 — CEO]}{강민수 CEO\rd{7mm}}}
 {\Bg{[검토자 — 전원]}{전원\rd{7mm}}}

%% ------------------------------------------------------------------ 1
\secbar{1. 세그먼트 정의와 크기}
\guide{초기 SAM(자동차 2·3차 30~300인 2,150개사)과 확장 SAM(11,800)을 적고, 선택 기준과 **회피하는 세그먼트**(예: 도면 클라우드 금지 1차사 계열)와 이유를 적는다.}
{\footnotesize
\begin{tabularx}{\linewidth}{|L{36mm}|Y|L{44mm}|}\hline
\hd{항목} & \hd{내용} & \hd{근거(정본 ID·인터뷰)} \\ \hline
\ifwbblank
\rd{10mm}초기 세그먼트(ICP) &  &  \\ \hline
\rd{10mm}크기(SAM) &  &  \\ \hline
\rd{10mm}선택 기준 3개 &  &  \\ \hline
\rd{10mm}회피 세그먼트와 이유 &  &  \\ \hline
\rd{10mm}확장 경로(언제·무엇이 되면) &  &  \\ \hline
\else
초기 세그먼트(ICP) & 자동차 2·3차 · 30~300인 · 라인 2~5 · MES 없음/엑셀 병행 · 1차사 포털 제출 의무 · 도면 클라우드 제약 없는 1차사 계열 & G-32, 27건, I-22·27 \\ \hline
크기(SAM) & 초기 SAM 2,150 중 클라우드 가능 1,247(58\%) [추가 설정]. 확장 SAM 11,800은 2027~ & G-32 \\ \hline
선택 기준 3개 & ① 종이 91\%(전기 시간 = 가치) ② 1차사 제출 의무(지불 논리가 외부에서) ③ IT 담당 없음 → 클라우드·폰만 & ㉒ 항목 4 \\ \hline
회피 세그먼트와 이유 & 도면 클라우드 금지 1차사 계열 903 — 온프렘은 운영·ISMS-P·GPU가 프리시드 범위 밖. 동보프레스는 보류 목록. 30인 미만(라인 과금 불성립) & I-27, X-03, A-07 \\ \hline
확장 경로 & ① 1,247 안에서 유료 5 → ② 1차사 채널로 온프렘을 채널 조건부 재검토(2027) → ③ 판정 라벨·마스터 구조가 같은 도메인부터 & ㉓ 항목 7 \\ \hline
\fi
\end{tabularx}}

%% ------------------------------------------------------------------ 2
\secbar{2. JTBD 진술문과 네 가지 힘}
\guide{'~한 상황에서, ~하기 위해, ~를 고용한다'. 힘 넷 — push(현 상황의 불만) / pull(새 해법의 매력) / anxiety(새 것에 대한 불안) / habit(현 방식의 관성)을 인터뷰 인용으로 채운다.}
{\footnotesize
\begin{tabularx}{\linewidth}{|L{30mm}|Y|L{40mm}|}\hline
\hd{항목} & \hd{내용} & \hd{인터뷰 근거} \\ \hline
\ifwbblank
\rd{10mm}JTBD 진술문 &  &  \\ \hline
\rd{10mm}push &  &  \\ \hline
\rd{10mm}pull &  &  \\ \hline
\rd{10mm}anxiety &  &  \\ \hline
\rd{10mm}habit &  &  \\ \hline
\else
JTBD — 검사원 & 라인이 돌아가는 동안, 나중에 나에게 책임을 물을 수 없도록 기록을 남기기 위해, 사진과 한마디로 끝나는 기록을 고용한다 & I-01·07·13·24 \\ \hline
JTBD — 품질팀장 & 부적합이 나왔을 때, 1차사에 늦지 않게 근거를 갖춰 통보하기 위해, 검사기록·사진·8D를 한 번에 만드는 것을 고용한다 & I-03·05·19·25 \\ \hline
push & '적는 게 일이지 검사가 일이 아니에요' · '하루만 빨랐어도' · 페널티 1,200만 & I-01·19 \\ \hline
pull & '사진은 찍어요, 그다음이 문제죠' · '그거 되면 제일 먼저 씁니다' & I-24·16 \\ \hline
anxiety & '다 해봤어요' · 판정이 기록으로 남는 것에 대한 침묵 · '도면이 나가면 잘립니다' · '우리 규모에 시스템은 사치' & I-02·27·14 — 막는 힘이 길다 \\ \hline
habit & 종이 일지는 라인 옆 · 전기는 퇴근 전 · 카톡 → 엑셀 개인 폴더 · 감사용 MES + 우리용 엑셀 & I-21 \\ \hline
\fi
\end{tabularx}}

%% ------------------------------------------------------------------ 3
\secbar{3. 도입 타임라인}
\guide{첫 생각 → 수동 탐색 → 능동 탐색 → 결정 → 소비 → 만족/불만. 각 단계에서 무슨 일이 있었나(인터뷰 사례 1건을 따라간다).}
{\footnotesize
\begin{tabularx}{\linewidth}{|L{30mm}|Y|L{50mm}|}\hline
\hd{단계} & \hd{무슨 일이 있었나(사례)} & \hd{우리가 개입할 지점} \\ \hline
\ifwbblank
\rd{9mm}첫 생각 &  &  \\ \hline
\rd{9mm}수동 탐색 &  &  \\ \hline
\rd{9mm}능동 탐색 &  &  \\ \hline
\rd{9mm}결정 &  &  \\ \hline
\rd{9mm}소비(도입 후) &  &  \\ \hline
\rd{9mm}만족/불만 &  &  \\ \hline
\else
첫 생각 & 2025 클레임 2건(4,800만) 후 '그때 누가 봤냐'를 못 댐 & 클레임 직후 대표에게 판정 로그(OP-04) \\ \hline
수동 탐색 & 2023 태블릿 → 6개월 후 종이 복귀 — '다 해봤다' & 실패 경험을 부정하지 않고 '입력이 아니라 촬영'으로 \\ \hline
능동 탐색 & 2026-03 인터뷰 요청 수락 — 정미영 팀장 & 문제 인터뷰 → WoZ 제안 \\ \hline
결정 & 2026-05 X-01 결과(47 → 19)를 본 뒤 LOI — 2026-09 시작, 구축비 별도 & 결과 숫자 + 조건표 \\ \hline
소비(도입 후) & 2026-09~ 3라인 파일럿 → D+240 첫 유료 & 온보딩 2주(A-08), 박선재의 회피 관찰 \\ \hline
만족/불만 & (미래) 정본 §3.7·§4.6 — 2027-08 사고 & A-05b가 열려 있다 \\ \hline
\fi
\end{tabularx}}

%% ------------------------------------------------------------------ 4
\secbar{4. 전략 커널 — 진단 / 추진 방침 / 일관된 행동}
\guide{Rumelt. 진단은 문제의 핵심 한 문장, 추진 방침은 그 문제를 다루는 접근 하나, 일관된 행동은 방침을 실행하는 3~5개 행동. 목표·비전을 쓰지 않는다.}
{\footnotesize
\begin{tabularx}{\linewidth}{|L{30mm}|Y|}\hline
\hd{요소} & \hd{내용} \\ \hline
\ifwbblank
\rd{11mm}진단 &  \\ \hline
\rd{11mm}추진 방침 &  \\ \hline
\rd{11mm}일관된 행동 1 &  \\ \hline
\rd{11mm}일관된 행동 2 &  \\ \hline
\rd{11mm}일관된 행동 3 &  \\ \hline
\else
진단 & 기록은 이미 만들어지고 있다(사진·종이·카톡) — 문제는 기록의 부재가 아니라 세 번 옮겨지며 근거를 잃는 것이고, 그 비용(31분·11.5h·'누가 봤냐')은 나뉘어 아무도 합산하지 않는다 \\ \hline
추진 방침 & 검사원의 행동을 바꾸지 않는다 — 이미 찍는 사진과 한마디에서 기록·근거·통보를 만들어 팀장에게 판다. 판정을 대신하지 않고 근거를 남긴다. 클라우드 가능 1,247에서 시작해 1차사 채널로 \\ \hline
일관된 행동 1 & 첫 6개월: 기록 자동 생성 + 부적합 패키지 → 포털 초안·8D 초안만. 판정 초안은 오전 라인 베타 \\ \hline
일관된 행동 2 & 라인 단위·좌석 무제한 — 검사원 수가 늘어도 비용이 안 늘어야 검사원이 쓴다 \\ \hline
일관된 행동 3 & 온프렘·커스텀 요구는 문서로 받고 문서로 보류(Won't, 재검토 조건) — 유료 5개사 전까지 · 행동 4: 판정 로그를 모든 기록에 \\ \hline
\fi
\end{tabularx}}

%% ------------------------------------------------------------------ 5
\secbar{5. 레퍼런스 클래스 원장 v0}
\guide{유사 버티컬 SaaS 12건 — 첫 유료까지 개월 / 첫 ARR 1억까지 개월 / 시드 규모. 우리 계획이 분포의 어디에 있는지 적는다(외부 시각).}
{\footnotesize
\begin{tabularx}{\linewidth}{|L{8mm}|Y|L{26mm}|L{26mm}|L{24mm}|L{40mm}|}\hline
\hd{\#} & \hd{회사·버티컬} & \hd{첫 유료까지(개월)} & \hd{ARR 1억까지(개월)} & \hd{시드 규모} & \hd{출처} \\ \hline
\ifwbblank
\rd{8mm} &  &  &  &  &  \\ \hline
\rd{8mm} &  &  &  &  &  \\ \hline
\rd{8mm} &  &  &  &  &  \\ \hline
\rd{8mm} &  &  &  &  &  \\ \hline
\rd{8mm} &  &  &  &  &  \\ \hline
\rd{8mm} &  &  &  &  &  \\ \hline
\rd{8mm} &  &  &  &  &  \\ \hline
\rd{8mm} &  &  &  &  &  \\ \hline
\else
1 & 국내 제조 검사기록 SaaS A & 9 & 22 & 시드 10억 & 투자사 \\ \hline
2 & 설비 예지보전 SaaS B & 14 & 30 & 시드 15억 & 공개 IR \\ \hline
3 & 식품 HACCP 기록 SaaS C & 6 & 18 & 프리시드 3억 & 멘토 \\ \hline
4 & 건설 안전점검 SaaS D & 8 & 21 & 시드 8억 & 공개 IR \\ \hline
5 & 일본 협력사 품질 포털 E & 12 & 26 & 시리즈A & 기사 \\ \hline
6 & 물류 현장 기록 SaaS F & 5 & 14 & 시드 6억 & 투자사 \\ \hline
7 & 병원 감염관리 기록 SaaS G & 11 & 28 & 시드 12억 & 공개 IR \\ \hline
8 & 농장 방제 기록 SaaS H & 7 & 24 & 프리시드 2억 & 멘토 \\ \hline
— & 중앙값 / 범위 (12건) & 9 / 4~14 & 21.5 / 14~30 &  &  \\ \hline
— & 딥게이지 계획 & 8 (D+240) — 중앙값 & 14 (D+420) — 상위 1/12 & 12억(예정) & ARR 계획은 분포 최상단 — IR에 쓰려면 근거(1차사 채널) \\ \hline
\fi
\end{tabularx}}

%% ------------------------------------------------------------------ 6
\secbar{6. 무엇을 하지 않는가}
\guide{전략은 하지 않는 것의 목록이다. 3~5개. 각 항목에 '하지 않는 대가'를 한 줄.}
\begin{fieldbox}
\ifwbblank
\lines{6}{9.5mm}
\else
1. 온프렘을 하지 않는다 — 대가: H사 계열 903개사와 동보프레스. 재검토: 유료 5개사 후 또는 1차사 채널 계약 시.
2. 자동차를 벗어나지 않는다 — 대가: 확장 SAM 11,800의 인바운드 거절. 재검토: 판정 라벨·마스터 구조가 같은 도메인부터, 2027~.
3. 판정을 대신하지 않는다(초안만) — 대가: 'AI 검사' 문구. 재검토: eval FN·FP 기준(㉗)과 검사원 수용(A-05b) 충족 시.
4. 공장의 문서·기준서 관리를 하지 않는다 — 대가: OP-06 요청 거절. 재검토: 1차사가 기준서 배포 채널로 지정할 때.
5. 컨설팅·구축 대행으로 돈을 벌지 않는다 — 대가: 초기 매출 느림. 구축비는 받되 ARR에 넣지 않는다(㉘).
\fi
\end{fieldbox}

\end{document}